% =============================================================================
% methods.tex — Complete Methods section (v0.3, 2026-09-03)
% Problem statement, base network model, GHG accounting, step demand,
% carbon policy modules A/B/C, quality-differentiated extension, solution
% strategy. Complements Model_Formulation_WI.tex (condensed formulation).
% =============================================================================

\section{Model Formulation}
\label{sec:model}

\begin{figure}[htbp]
\centering
\includegraphics[width=\textwidth]{figures/fig0_superstructure.png}
\caption{\small Biochar supply chain superstructure and policy architecture.
(I) Material and carbon flows: 13 feedstocks gathered at 72 county nodes
(representative nodes $n_1$, $n_2$), T1 dehydration, T2 slow pyrolysis at
300/500$^\circ$C, biochar delivered to the price-responsive step demand
($f^{H},f^{M},f^{L}$) and to the non-soil sink ($q^{\mathrm{sink}}$), and the
carbon-crediting layer: 500$^\circ$C chars (H/C 0.40--0.52) are
credit-eligible while torrefaction-like 300$^\circ$C chars (H/C 0.9--1.4) are
not. (II) The three policy architectures differ in accounting basis,
eligibility rule, and instrument, so they are shown as separate modules rather
than a single flow: (A) a throughput-scaled baseline-and-credit whose credit
$R=B-E^{+}+|S|$ is paid per dry tonne with no eligibility gate; (B) a
gross-emission cap $E^{+}\leq\mathrm{CAP}$ whose shadow price $-\mu$ is the
marginal abatement cost; and (C) a tiered tax combined with an H/C-gated
VM0044-inspired credit. In (II), italic row labels give the material, carbon,
and cash flows of each module, dashed rules mark the constraint boundary, and
the symbol key defines all notation.}
\label{fig:super}
\end{figure}

\subsection{Problem statement}

We consider a regional biochar supply chain operator (a coordinating entity,
per the market-clearing paradigm of \cite{sampat2019,ma2023}) that must decide where to
build biomass drying and pyrolysis facilities, which feedstocks to procure from
which counties, and how to distribute biochar to demand locations, in order to
maximize the system-wide economic surplus (Fig.~\ref{fig:super}). The network
consists of $n \in
\mathcal{N}$ county nodes ($|\mathcal{N}|=72$ for Wisconsin), $f \in \mathcal{F}$
feedstocks ($|\mathcal{F}|=13$: corn stover, soybean straw, wheat straw, oats
straw, barley straw, switchgrass, miscanthus, poplar, willow, logging residues,
small-diameter trees, mill residues, and other forest waste), and a two-stage
technology superstructure: T1 dehydration followed by T2 slow pyrolysis at
either 300 or 500~$^\circ$C. The model jointly determines (i) integer facility
locations and scale tiers, (ii) continuous material flows, and (iii) the
equilibrium quantities served on a price-responsive biochar demand curve. Under
carbon policy, the operator additionally earns removal credits or faces
emission charges, which enter the objective as described in
Section~\ref{sec:policy}.

\subsection{Base network model}

The base model is a multi-product supply chain network design formulation; the
complete statement (with sets, parameters, and the full constraint list) is
given in the SI (Table~S5). Throughout, sets are calligraphic, decision
variables are lowercase, parameters are Greek or barred, and $\delta_{ij}$
denotes the great-circle distance between nodes $i$ and $j$. The operator
maximizes the coordinated-system surplus
\begin{subequations}
\label{eq:base}
\begin{align}
\maxi_{z,x,f,\mathrm{dem},\mathrm{sup}}\quad
\Pi & = \underbrace{\textstyle\sum_{d\in\mathcal{D}} \lambda^{\mathrm{dem}}_{d}\,\mathrm{dem}_{d}}_{\text{revenue}}
- \underbrace{\textstyle\sum_{s\in\mathcal{S}} \lambda^{\mathrm{sup}}_{s}\,\mathrm{sup}_{s}}_{\text{farmgate}}
- \underbrace{\textstyle\sum_{n,t} \omega_{t}\,(-x_{n,\mathrm{ref}(t),t})}_{\text{OPEX}}
\notag\\
&\quad - \underbrace{\textstyle\sum_{(i,j,p)\in\mathcal{A}}(\alpha^{\mathrm{var}}_{p}\delta_{ij}+\alpha^{\mathrm{fix}}_{p})f_{ijp}}_{\text{transport}}
- \underbrace{\mathrm{AF}\textstyle\sum_{n,t,k}\kappa_{t,k}\,z_{ntk}}_{\text{annualized CAPEX}}
\label{eq:base_obj}\\
\st & \quad s_{np} + p_{np} + \textstyle\sum_{j:(j,n,p)\in\mathcal{A}} f_{jnp} =
d_{np} + \textstyle\sum_{j:(n,j,p)\in\mathcal{A}} f_{njp},
\quad (\pi_{np}) \quad n\in\mathcal{N},\ p\in\mathcal{P} \label{eq:base_balance}\\
& \quad x_{n,p,t} = \gamma_{t,p}\,x_{n,\mathrm{ref}(t),t}, \quad n\in\mathcal{N},\ t\in\mathcal{T},\ p\in\mathcal{P} \label{eq:base_yield}\\
& \quad -x_{n,\mathrm{ref}(t),t} \le \textstyle\sum_{k} \bar{q}_{t,k}\,z_{ntk}, \quad n\in\mathcal{N},\ t\in\mathcal{T} \label{eq:base_cap}\\
& \quad z_{ntk}\in\{0,\dots,5\}, \quad 0\le f_{ijp}\le \bar{f}_{ijp}, \quad \mathrm{dem}_d\le\bar{d}_d, \quad \mathrm{sup}_s\le\bar{s}_s \label{eq:base_bounds},
\end{align}
\end{subequations}
where $\mathcal{P}$ collects all products (wet biomass, dry biomass, biochar,
carbon credits), $s_{np}$/$d_{np}$ aggregate the supply and demand bids at node
$n$, and $\gamma_{t,p}$ is the fixed-yield technology transformation.
Constraint \eqref{eq:base_obj} captures demand-side surplus from the bid stack
minus farmgate, operating, transport, and annualized capital costs.
Constraint \eqref{eq:base_balance} is the node mass balance for product $p$;
its dual variable $\pi_{np}$ is the spatial inherent value of product $p$ at
node $n$, the marginal cost of delivering one more unit of $p$ at $n$, which
the market-clearing paradigm interprets as the equilibrium price
\cite{sampat2019,ma2023}. Constraint \eqref{eq:base_yield} enforces the
fixed-yield technology transformations (dehydration water removal, pyrolysis
mass yields, credit co-production), \eqref{eq:base_cap} links throughput to the
integer number of facilities installed at each scale tier $k$, and
\eqref{eq:base_bounds} states integrality and the bid capacities. Transport
radius limits (wet $\leq$150~km, dry $\leq$250~km, biochar $\leq$400~km) are
imposed through $\bar{f}_{ijp}$, consistent with biomass-logistics cost
estimates under which hauling low-density feedstocks becomes prohibitive
beyond roughly 80--160~km \cite{searcy2007}. The annualization factor is
$\mathrm{AF}=0.1175$ ($i{=}10\%$, $n{=}20$~yr); all parameter values and data
sources are tabulated in the SI.

\subsection{Greenhouse gas accounting}

Three GHG components enter the policy modules, all in kg CO$_2$e/yr. Gross
emissions $E^{+}(x)$ comprise process emissions from collection, drying, and
pyrolysis plus transport:
\begin{equation}
E^{+}(x)=\sum_{n,t} e^{\mathrm{proc}}_{t}\bigl(-x_{n,\mathrm{ref}(t),t}\bigr)
+ e^{\mathrm{tr}}\sum_{(i,j,p)\in\mathcal{A}}\delta_{ij}f_{ijp},
\end{equation}
with $e^{\mathrm{proc}}_t$ the technology-specific process intensity per
tonne of the technology's reference product (31--35~kg/t wet, i.e.,
35--51~kg/t dry, for T1, 113--122~kg/t dry for T2 at 500/300$^\circ$C after
the syngas electricity credit)
and $e^{\mathrm{tr}}=0.0902$~kg/t-km. Carbon removal $S(x)$ is defined as a
positive physical quantity: sequestration and N$_2$O suppression per
tonne of biochar applied:
\begin{equation}
S(x)=\sum_{n,\,t\in\mathcal{T}^{\mathrm{pyr}}}
\bigl(r^{\mathrm{seq}}_{t}+r^{\mathrm{N}_2\mathrm{O}}_{t}\bigr)\,x_{n,\mathrm{BC},t},
\qquad
r^{\mathrm{seq}}_{t}=\rho_t\,C_t\,\frac{44}{12}\times 10^{3}\ \text{kg/t},
\end{equation}
where $C_t$ is the biochar carbon content, $\rho_t$ a permanence factor
($\rho_{300}=0.65$, $\rho_{500}=0.80$; VM0044 v1.2 Table~3, drawn from IPCC
2019 Table~4AP.2; the 300$^\circ$C value uses the lowest tabulated band,
350--450$^\circ$C, and the permanence of sub-350$^\circ$C chars may be lower),
and $r^{\mathrm{N}_2\mathrm{O}}_{t}=52$~kg/t from the
meta-analysis of Cayuela et al.~\cite{cayuela2014}. The direct atmospheric flux of the system is therefore
\begin{equation}
N(x)=E^{+}(x)-S(x),
\end{equation}
which is negative (net removal) at the baseline equilibrium. This flux
excludes the decomposition baseline, the emissions avoided by diverting
biomass from natural decay,
\begin{equation}
B(x)=\sum_{f\in\mathcal{F}} b_f \sum_{n}\bigl(-x_{n,f,\mathrm{T1}_f}\bigr),
\qquad b_f = 0.9\,C_f^{\mathrm{dry}}\,\frac{44}{12}\times 10^{3}\ \text{kg/t dry},
\end{equation}
with $b_f \approx 1350$--$1650$~kg/t dry. The decay fraction $0.9$ is a
scenario assumption (immediate, complete decay of uncollected residues, an
upper bound on the avoided-emission credit); a sensitivity halving it to 0.5 is
reported in Section~\ref{sec:sensitivity}. Throughout the paper, reported
``net GHG'' figures are the direct flux $N(x)$, not abatement relative to
the decay counterfactual; $B(x)$ enters only the credit calculation of
Paradigm~A below.

\subsection{Price-responsive demand and market clearing}

Biochar demand follows a quality-tiered step function (bid stack) at each
county, allocated in proportion to farmland area:
\begin{equation}
\lambda^{\mathrm{dem}}(Q)=\begin{cases}
\$800/\mathrm{t} & Q\le 0.8~\mathrm{Mt}\ \text{(CDR premium: BC500 only in v2b)}\\
\$500/\mathrm{t} & 0.8<Q\le 2.3~\mathrm{Mt}\ \text{(quality ag / industrial media)}\\
\$250/\mathrm{t} & 2.3<Q\le 4.8~\mathrm{Mt}\ \text{(bulk agricultural)}\\
\$30/\mathrm{t}  & \text{sink (construction, filtration)}
\end{cases}
\end{equation}
The three segment prices are anchored to production costs of \$150--400/t
\cite{roberts2010,shabangu2014} plus the value of the co-produced removal
credit (\$100--250/tCO$_2$e \cite{cdrfyi2024,verraVM0044}; $\approx$1.6~tCO$_2$e
per tonne of 500$^\circ$C char, i.e.\ \$160--400/t of biochar) and a quality
premium for durable, H/C-gated product: the \$800/t H tier is durable biochar
sold to national corporate CDR buyers, the \$500/t M tier quality agricultural
and industrial-media grades, and the \$250/t L tier bulk soil application.
This segmentation follows the soil-amendment, industrial-media, and
carbon-removal end uses that structure the biochar market, within the wide
price dispersion reported in current market surveys \cite{ibi2025}. The three soil
segments sum to a finite aggregate capacity of $0.8+1.5+2.5 = 4.8$~Mt/yr,
which is the ceiling on price-responsive biochar sales; the sink is unbounded
(open capacity) and absorbs any residual at the low \$30/t bid. All three
capacities are allocated across counties in proportion to farmland area, a
smooth spatial proxy rather than observed county-level demand; because every
county then carries demand at the same bid, the binding constraint is the
aggregate demand quantity and the equilibrium spatial pattern is
feedstock-driven rather than demand-driven (SI, demand data pack).
Because demand is
price-responsive, the equilibrium quantity is endogenous; in the base model
(v2) biochar is pooled, while the quality-differentiated extension (v2b,
Section~\ref{sec:v2b}) splits biochar into BC300 and BC500 products with the
CDR premium segment restricted to BC500.

\emph{Price concepts.} Four distinct price-like objects appear throughout the
paper, and they are not the same quantity. (i) \emph{Demand bid}
$\lambda^{\mathrm{dem}}_d$: the buyer's willingness-to-pay anchor for one tonne
of biochar in segment $d$, not an equilibrium price. (ii) \emph{Nodal inherent
value} $\pi_{np}$: the dual of the node mass balance, the conditional marginal
value of one more tonne of product $p$ at node $n$ under the fixed facility
layout. (iii) \emph{Shadow carbon price} (conditional fixed-layout marginal
abatement cost): the negative of the cap dual, $-\mu_{\mathrm{CAP}}$, the
marginal opportunity cost of tightening the emission cap under the fixed
layout; it is a coordinating-entity shadow price, not a price from a
decentralized allowance market (Section~\ref{sec:limitations}).
(iv) \emph{Administered market benchmarks}: observed RGGI/EU ETS allowance and
voluntary removal-credit prices, used only as external comparisons.

\subsection{Carbon policy modules}
\label{sec:policy}

\subsubsection{Paradigm A: baseline-and-credit}

Credits accrue for the net emission reduction relative to the decomposition
baseline, and because the baseline scales with throughput, the credit mechanism
rewards capacity expansion, the property a linear tax lacks:
\begin{equation}
R(x)=B(x)-E^{+}(x)+S(x)\ \ge\ 0,
\qquad \max\ \Pi+\frac{p_c}{1000}R(x),
\qquad p_c\in\{25,50,100,150,200\}\ \mathrm{\$/tCO}_2\mathrm{e}.
\end{equation}
Implementation: carbon credit (CC) is a co-produced product with per-tonne-dry
rates $\gamma_{t,28}=(b_f-e^{\mathrm{proc}}_{t}+(r^{\mathrm{seq}}_t+r^{\mathrm{N}_2\mathrm{O}}_t)\gamma_{t,\mathrm{BC}})/1000$,
sold at the policy price $p_c$. Credits are computed in LCFS style for
every feedstock--temperature pair, without an eligibility gate; the
VM0044 H/C gate is applied only in Paradigm~C2, so the two paradigms differ in
both credit basis (avoided-baseline versus sequestered-carbon) and eligibility
(Fig.~\ref{fig:creditbasis}). Paradigm~A is therefore an idealized
throughput-scaled credit rather than a specific existing methodology; it is
contrasted with the sequestration-only, H/C-gated VM0044 basis of Paradigm~C2.

\subsubsection{Paradigm B: cap-and-trade with shadow carbon price}

B1 imposes a gross-emission cap and sweeps it to trace the system's marginal
abatement cost (MAC) curve $\Pi^{*}(\mathrm{CAP})$:
\begin{equation}
\max\ \Pi \quad \text{s.t.}\quad E^{+}(x)\le \mathrm{CAP},
\qquad \mathrm{MAC}=\frac{\partial \Pi^{*}}{\partial\,\mathrm{CAP}}
\end{equation}
The shadow-price interpretation is formalized as follows.
\begin{proposition}
\label{prop:mac_dual}
For the fixed-layout LP \eqref{eq:base} with an added gross-emission cap
$E^{+}(x)\le\mathrm{CAP}$, let $\mu_{\mathrm{CAP}}$ be the optimal dual
variable of the cap constraint under Gurobi's sign convention for a $\le$
constraint in a maximization. Then
$\mu_{\mathrm{CAP}}=-\partial\Pi^{*}/\partial\,\mathrm{CAP}$: the reported dual
is the negative of the marginal abatement cost, so the shadow carbon price
(conditional on the fixed layout) is $-\mu_{\mathrm{CAP}}$.
\end{proposition}
\begin{proof}
The identity is the standard parametric-LP sensitivity theorem applied to the
cap's right-hand side; the full argument is in the SI (Section~A.9).
\end{proof}
The proposition is verified numerically: the reported dual matches the
finite-difference MAC to numerical tolerance ($-374$ vs.\ \$351--391\$/t over
the first binding intervals in the near-term scenario). We benchmark the
resulting MAC against 2024--2025 RGGI
(\$17--25/t) and EU ETS (\texteuro 50--90/t, $\approx$\$55--100/t) allowance
prices and the
voluntary removal-credit band (\$100--250/t, VM0044/Puro). B2 replaces the
constraint by a net-flux cap
$E^{+}(x)-S(x)\le\mathrm{CAP}^{\mathrm{net}}$. At every tested cap
($\mathrm{CAP}^{\mathrm{net}}\ge-1{,}000$~kt, all looser than the baseline
flux of $-1{,}595$~kt) the constraint does not bind; a frontier sweep that
tightens the cap below the baseline flux (free-facility MIP, SI Fig.~S2)
locates the first binding point and traces the response, which is reported in
Section~\ref{sec:results}.

\subsubsection{Which resource binds: demand-bound versus supply-bound regimes}

Whether a policy instrument changes the technology mix is decided by
which resource binds. Define, for each feedstock $f$, the dry-biomass
mass yields of the 300$^\circ$C and 500$^\circ$C pyrolysis routes
$\gamma^3_f$ and $\gamma^5_f$ ($\gamma^3_f>\gamma^5_f$), the credit rates
$\rho^3_f, \rho^5_f$ (tCC/t dry, zero for ineligible 300$^\circ$C chars under
VM0044), the per-tonne dry cost $\omega$ of the full chain (farmgate, OPEX,
transport, annualized CAPEX), and a biochar market price $P$. The two
candidate ``regimes'' correspond to which constraint is active in the
fixed-facility LP:
\begin{itemize}
\item \textbf{Demand-bound} (feedstock abundant, e.g.\ mature-market): the
supply of biochar is limited by demand, so the relevant objective margin is
per tonne of biochar served. The 500$^\circ$C margin premium per
served tonne is
\begin{equation}
\Delta^{\mathrm{dem}}_{f}(p_c)=
\Bigl(\tfrac{\rho^5_f}{\gamma^5_f}-\tfrac{\rho^3_f}{\gamma^3_f}\Bigr)p_c
-\omega\Bigl(\tfrac{1}{\gamma^5_f}-\tfrac{1}{\gamma^3_f}\Bigr),
\end{equation}
independent of the market price $P$: a lower yield multiplies credit
revenue per tonne served.
\item \textbf{Supply-bound} (feedstock scarce, e.g.\ near-term): the supply of
biochar is limited by feedstock, so the relevant margin is per tonne of dry
biomass processed. The 500$^\circ$C premium per dry tonne is
\begin{equation}
\Delta^{\mathrm{sup}}_{f}(p_c)=
\bigl(\gamma^3_f-\gamma^5_f\bigr)P
-\bigl(\rho^5_f-\rho^3_f\bigr)p_c,
\end{equation}
independent of the chain cost $\omega$: the yield loss is now a revenue loss
($P$ scales it) and the credit-rate gain is compared per dry tonne.
\end{itemize}
\begin{proposition}
\label{prop:regime}
For a fixed feedstock $f$ and biochar market price $P$:
(i) in the demand-bound regime, the 500$^\circ$C route becomes optimal at
the unique crossing price
\begin{equation}
\bar{p}^{\mathrm{dem}}_f = \frac{\omega\,(1/\gamma^5_f-1/\gamma^3_f)}
{\rho^5_f/\gamma^5_f-\rho^3_f/\gamma^3_f},
\end{equation}
which does not depend on $P$ and is strictly positive whenever
$\rho^5_f/\gamma^5_f>\rho^3_f/\gamma^3_f$; and
(ii) in the supply-bound regime, the 500$^\circ$C route becomes optimal at
\begin{equation}
\bar{p}^{\mathrm{sup}}_f = \frac{(\gamma^3_f-\gamma^5_f)\,P}
{\rho^5_f-\rho^3_f},
\end{equation}
which is increasing in $P$ and bounded by the demand-segment price.
Moreover, for the LCFS-style baseline-and-credit rates of Paradigm~A
(Table~S2: feedstock averages $\bar{\gamma}^{3}=0.352$,
$\bar{\gamma}^{5}=0.278$ t BC/t dry, $\bar{\rho}^{3}=1.771$,
$\bar{\rho}^{5}=1.849$ tCC/t dry (the per-dry-tonne rates, whose
$\approx 0.077$ tCC/t dry premium is the credit-rate advantage of
500$^\circ$C)) and the pooled chain cost $\omega\approx\$130$/t dry,
$\bar{p}^{\mathrm{dem}}\approx\$61$/tCO$_2$e for every demand segment and
$\bar{p}^{\mathrm{sup}}\approx\$239$/tCO$_2$e at $P=\$250$/t.
\end{proposition}
\begin{proof}
Equating the per-tonne margins of the two routes in each regime and solving
for $p_c$ yields the stated crossing prices; the full argument is in the SI
(Section~A.10).
\end{proof} the policy--resource interaction observed in
Section~\ref{sec:results}: baseline-and-credit never
upgrades the technology mix in the near-term (feedstock-scarce) scenario up to
$p_c\approx\$239$/t, whereas the mature-market (demand-bound) scenario
switches at $p_c\approx\$61$/t for every demand segment; the quality
tier of the premium segment ($\$800$~t$^{-1}$, BC500-only) is unaffected,
since it is already served by 500$^\circ$C in both regimes. The proposition
also explains why Paradigm~C2's VM0044 gate is so much more powerful than a
flat credit price: the gate sets $\rho^3_f=0$, maximizing the numerator in
$\bar{p}^{\mathrm{sup}}_f$, i.e.\ it prices durability rather than tonnes.

\subsubsection{Paradigm C: tiered tax + VM0044 eligibility + allowance allocation}

C1 taxes gross emissions on a convex piecewise-linear schedule with tiers
calibrated inside the feasible emission range ($T_1=0.4E^{+}_0$, $T_2=0.7E^{+}_0$;
rates \$25/50/100 and \$50/100/200 per tCO$_2$e); the flat rate set
(25/25/25) isolates the linear-tax case. To rule out the possibility that the
fixed-layout approximation drives the tax results, C1 is also re-solved as a
free-facility MIP (C1-free). C2 adds removal crediting operationalized through
the molar H/C~$\le$~0.7 eligibility gate as an H/C-gated VM0044-inspired
proxy:
only biochar with molar H/C~$\le$~0.7 generates credits. Measured 300$^\circ$C
chars are torrefaction-like with H/C $\approx$0.9--1.4
\cite{rafiq2016,ippolito2020}, so no 300$^\circ$C biochar qualifies and
its credit rate is zero for every feedstock; 500$^\circ$C chars have H/C
$\approx$0.40--0.52 \cite{rafiq2016,enders2012,ippolito2020} and qualify for
all feedstocks, turning carbon crediting into a technology-selection
instrument. Other VM0044 requirements (additionality
demonstration, leakage accounting, monitoring and verification) are outside the
model scope and their transaction costs are not modeled. C3 compares allowance
allocation rules: grandfathering ($A=E^{+}_0$, fixed) versus benchmarking
($A=\beta q$ with $\beta=E^{+}_0/q_0$), with allowance trades priced at
$p_{\mathrm{allow}}\in\{30,80\}$~\$/tCO$_2$e.

\subsubsection{Why a linear emission tax does not change a demand-constrained system}

The observed null response of the tiered tax on dispatch (Section~\ref{sec:policyC})
is not a modeling artifact but a structural property of the demand-constrained
regime. Formally:
\begin{proposition}
\label{prop:taxinv}
Let $x^{*}$ be an optimal solution of the fixed-layout LP \eqref{eq:base} and
suppose the biochar mass-balance duals (inherent values) satisfy
$\pi^{\mathrm{BC}}_{n} \le \lambda^{\mathrm{dem}}_{d(n)}$ for every served
demand bid $d(n)$ and $\pi^{\mathrm{BC}}_{n} \ge
\lambda^{\mathrm{dem}}_{d}$ for every unserved bid, both with strict
inequalities (the marginal served bid earns a strictly positive margin and
the next, unserved, bid is strictly unprofitable; this is the demand-constrained
regime with a strict gap at the margin). Then for any tax rate
$\tau \in [0, \tau_{\max})$ with
\begin{equation}
\label{eq:taumax}
\tau_{\max} =
\min_{q\in\mathcal{Q}(x^{*})}
\frac{\lambda^{\mathrm{dem}}_{d(q)}-\pi^{\mathrm{BC}}_{n(q)}}{e^{+}_q},
\end{equation}
where $\mathcal{Q}(x^{*})$ is the set of biochar supply paths that are
marginally profitable at $x^{*}$
($\lambda^{\mathrm{dem}}_{d(q)}>\pi^{\mathrm{BC}}_{n(q)}$) and $e^{+}_q$ is
the gross-emission rate of path $q$ (kg CO$_2$e per tonne of biochar), the
same $x^{*}$ remains optimal in the taxed problem
$\max\ \Pi(x)-\tau\,E^{+}(x)$; the tax reduces only surplus,
$\Pi_{\tau}(x^{*})=\Pi(x^{*})-\tau\,E^{+}(x^{*})$, and leaves the material
flows, technology mix, and equilibrium quantities unchanged.
\end{proposition}
\begin{proof}
The tax perturbs only the objective, so the optimal basis is unchanged while
the reduced-cost signs of the active margin paths persist; the full argument is
in the SI (Section~A.11).
\end{proof}
Intuitively, the tax rescales the objective by a constant per-unit-of-activity
penalty that is uniform across the supply paths whose activity is fixed by the
marginal demand bid; as long as the (gross) emission intensity of those paths
is bounded relative to their surplus, the pre-tax ordering of the bid stack is
preserved. In our data, the marginal price gap is
$\lambda^{\mathrm{dem}}-\pi^{\mathrm{BC}} \approx \$155$/t
biochar (the \$500/t M bid minus the largest served-node inherent value,
\$344.85/t). Evaluating \eqref{eq:taumax} with the route-specific marginal
intensities of the LP basis (SI Section~A.11 reports every served demand row
individually), the binding rows are the M-segment rows served by the
corn-stover 300$^\circ$C chain with $e^{+}_q\approx0.48$~t CO$_2$e/t biochar,
so $\tau_{\max}\approx\$318$/t CO$_2$e; the H-segment rows are far looser
($\approx$\$1{,}100/t). In the mature-market scenario the marginal gap is
wider and $\tau_{\max}\approx\$464$/t, so the neutrality result strengthens
with feedstock abundance. The route-specific threshold exceeds the top tested
tier of
\$200/t; the tax's effect within the tested range is confined to surplus
reallocation. The neutrality itself is a structural property of any
demand-constrained LP; the contribution here is the quantified boundary
($\tau_{\max}\approx\$318$/t for this system) and the comparison it enables
with observed carbon prices.

\subsection{Quality-differentiated extension (v2b)}
\label{sec:v2b}

The base model pools biochar into one product. The extension splits it into
BC300 (product 27, from 300$^\circ$C pyrolysis, torrefaction-like H/C
$\approx$0.9--1.4, credit-ineligible) and BC500
(product 28, from 500$^\circ$C, H/C~$\approx$~0.40--0.52, credit-eligible), with
the CDR premium demand
segment ($\$800/t$) restricted to BC500. Yields follow the V9 technology table,
cross-checked against \cite{roberts2010} and \cite{rafiq2016}:
$\gamma^{\mathrm{BC300}}=0.34$--$0.38$ and $\gamma^{\mathrm{BC500}}=0.29$--$0.32$,
with carbon contents 41--50\% and 50--61\%, respectively (the upper 500$^\circ$C
carbon contents are conservative relative to the 62--65\% published range,
Table~S2). This makes the
pyrolysis temperature a genuine economic decision: 500$^\circ$C sacrifices
yield to access the premium market and higher removal credits.

\subsection{Solution strategy}

Facility-location decisions are integer, so Steps 1 (MIP) and the policy sweeps
use Gurobi 13 with 0.5--1\% optimality gap targets, 600--3{,}600~s time limits, and
feasibility-focused branching (MIPFocus=1) with 60~s of zero-objective
heuristics; successive policy prices are warm-started from the previous
facility layout. All dual-based analyses (inherent values, GHG shadow prices, shadow
carbon prices) are performed on LPs with the facility layout fixed
and solved with dual simplex, so that reported duals are vertex
shadow prices; where degeneracy is detected we report the dual range instead
of a point value (SI Section~A.9), and interior-point (barrier) duals of the
cap constraint are not
vertex shadow prices on this LP and are not used. The GHG LP (Step 3) pins demand quantities to the
Step-2 equilibrium, so its duals are marginal GHG values at the market
equilibrium rather than at the GHG-minimal dispatch. The near-term scenario
(6.97~Mt dry biomass) serves as the baseline and the mature-market medium
scenario (13.53~Mt dry) as sensitivity. Full reproducibility scripts and data
are provided in the SI.

\begin{assumption}
\label{ass:duality}
For the fixed-layout LPs, strong duality holds: the reported primal and dual
objective values coincide. Dual solutions need not be unique at degenerate
vertices, so we do not assume uniqueness; throughout, the reported shadow
prices are verified against finite-difference perturbations, and dual ranges
(SI Section~A.9) are reported wherever degeneracy is detected.
\end{assumption}
